\pagenumbering{Roman}
\thispagestyle{empty}
\begin{center}
{\fontsize{16pt}{20pt}\selectfont\bfseries 无线电干扰源的主动定位与搜索清除\par}
\vspace{5pt}
\end{center}
% ===== 摘要初稿：目前概括已写的一、二、三问，可直接在此修改 =====
\vspace{12pt}
\begin{center}
{\fontsize{14pt}{17pt}\selectfont\bfseries 摘\quad 要\par}
\end{center}
\vspace{4pt}

针对测角误差、有限接收范围与多类动作耗时共同约束下的无线电干扰源定位清除任务，本文以源位置的可行区域描述观测信息，建立从单源交会定位到多源自动搜索的数学模型，在保证可靠清除的前提下统筹检测位置、频道选择与任务顺序。

\textbf{针对问题一，}将各次示向度的角度约束转化为半平面交，构造凸多边形定位区域；利用凸多边形顶点及旋转卡壳计算区域直径，并通过反例说明，以最远点对连线为直径的圆不一定覆盖整个定位区域，从而区分区域直径与覆盖圆的几何含义。

\textbf{针对问题二，}结合首测信息与接收半径下界，构造保证再次接收信号的第二检测点候选区域。依据定位区域长轴生成有限探点，综合名义观测后的区域收缩与移动耗时进行评分，给出可执行的主动选点策略，并通过已有双测算例展示定位区域的收缩过程。

\textbf{针对问题三，}将发现覆盖与已知源的定位清除统一为滚动决策任务，采用位掩码状态压缩及有预算的 A* 搜索安排任务顺序，结合共享测量、静默信息删测与可靠清除判据完成多源搜索。以前瞻策略为基线，与 PPO 强化学习及几何法开展同场景比较，并通过消融实验和轨迹对照分析主要策略的作用。

\textbf{针对问题四，}将定向源的辐射范围写成过源位置的半平面，指出无信号反馈在本问不再对应距离信息；据此把发现覆盖条件改为“候选源位置落在其 $1000$ m 内检测站的凸包中”，用原点、内环 $8$ 站与场外外环 $12$ 站共 $21$ 个检测站逐格验证该条件。对已发现的源采用成对检测点，使两点同时无信号时仍能截断定位区域，并用确定在接收范围内的无信号点联合有信号点限制发射方向和位置。任务安排沿用滚动规划，处理单个源时用与历史反馈相符的假想源模拟比较候选动作，并用离线训练的回归树判断是否在途中停下补测。

在既有 $256$ 个随机场景中，首版状态搜索均完成全部干扰源清除，平均整局计费用时为 $3081.41\,\mathrm{s}$，较前瞻基线降低 $8.34\%$。在按附件接口复原的本地模拟器上，问题四的策略在 $1000$ 个随机混合源场景中全部完成清除，平均定位清除时间为 $449.9\,\mathrm{s}$/源，比只用几何规则的方案降低 $1.0\%$，对贴着边界朝外发射的压力场景降低 $9.2\%$。结果说明，利用几何约束压缩决策空间并联合规划检测与清除任务，能够在保证不漏源、不误清的前提下改善平均任务效率。

\vspace{12pt}
\noindent\textbf{关键词：}交会定位；主动测点选择；状态压缩；A* 搜索；强化学习；定向干扰源

